% !TeX root = ../main.tex
\section{Discussions and Conclusions}
\subsection{Discussion of Results}
This thesis set out to examine whether computer vision models for soccer analysis could be trained, compressed, and deployed effectively on embedded hardware without sacrificing too much analytical performance. The results indicate that this is achievable, but only when compression is applied selectively and interpreted in the context of the target task. Across both object detection and pitch keypoint localisation, the most useful deployment candidates were not the most aggressively compressed models, but the configurations that preserved enough accuracy while still reducing latency, power draw, and thermal load to make edge deployment practical.

A clear theme throughout the results is that hardware acceleration should be treated as the baseline optimisation step rather than as an optional enhancement. Figure~\ref{fig:hardware_acceleration_results} showed that TensorRT-based FP32 acceleration reduced latency substantially while preserving almost identical mAP$_{50\text{--}95}$. This is important because it means a considerable portion of the deployment benefit can be obtained before any accuracy-sensitive compression is introduced. In practical terms, the acceleration stage improves real-time viability without changing the relative ranking of the models, so model selection can still be based primarily on the original accuracy--latency balance.

Quantisation emerged as the strongest method for improving embedded efficiency overall. The latency, power, and temperature results consistently showed that INT8 deployment produced the largest reductions in computational cost. However, Figure~\ref{fig:quantisation_results} and the summary tables also showed that these gains came with a more explicit accuracy penalty than hardware acceleration alone. This trade-off is especially important for edge sports analytics, where improved throughput is valuable only if the extracted statistics remain trustworthy. The quantisation results therefore suggest that INT8 deployment is most appropriate when strict runtime or energy limits dominate, whereas FP32-accelerated models remain preferable when analytical fidelity is the primary concern.

One of the most important findings is that compression tolerance depends strongly on the vision task. The object detection models were consistently more robust than the pitch keypoint models under pruning and combined compression paths. Figure~\ref{fig:pruning_results} suggests that the object models retained enough representational redundancy for structured pruning to reduce latency without catastrophic accuracy loss, particularly at larger scales. In contrast, the pose models degraded much more quickly as pruning strength increased. A plausible interpretation is that pitch keypoint localisation depends on finer spatial structure and more fragile feature relationships than player and ball detection, making it less tolerant to parameter removal and reduced numerical precision. This task-dependent sensitivity means that a single compression recipe cannot be assumed to generalise across all stages of a sports analytics pipeline.

Knowledge distillation occupied a different role from the other techniques. As shown in Figure~\ref{fig:knowledge_distillation_results}, distillation did not materially change latency on its own, but it often helped smaller models retain or slightly improve analytical performance. This makes distillation less valuable as a direct deployment optimisation and more valuable as a supporting mechanism that stabilises compressed students. The combined-path analysis in Figure~\ref{fig:combination_results} reinforces this interpretation, since the strongest deployment candidates were typically balanced intermediate configurations in which acceleration and moderate compression were paired with accuracy-preserving steps such as distillation. In other words, the results support a staged view of optimisation: first accelerate, then compress cautiously, and use distillation where accuracy recovery is needed.

The hardware measurements also add an important practical layer to the discussion. Improvements in latency did not always translate proportionally to reductions in GPU or CPU usage, and quantisation in particular sometimes increased CPU usage for larger models. This suggests that some deployment benefits are achieved by shifting the system bottleneck rather than eliminating it entirely. For an embedded platform, that distinction matters. A model that is faster but transfers more overhead to the host may still be useful, but it changes the design constraints for the rest of the pipeline. The power and temperature results therefore strengthen the case for quantisation as a deployment tool, since they show that its value is not limited to faster inference alone, but extends to lower sustained operating cost and improved thermal behaviour.

The full-stack replay experiments show that these model-level findings also translate into system-level usefulness. The best full-stack configuration achieved a Pass mAP$_{1\text{--}5}$ of 0.697, 13.02 FPS, and 39.61\% GPU utilisation, demonstrating that embedded soccer analytics can be carried out on the Jetson Orin Nano without saturating the device. At the same time, these results also show that the original real-time target of roughly 20 FPS was not fully reached by the final end-to-end pipeline. This is an important nuance. The thesis still demonstrates practical viability, especially for statistics such as possession and pass analysis that evolve over multiple frames, but it also shows that model compression alone is not sufficient to close the remaining gap to higher-throughput real-time operation. System-level optimisations in scheduling, detection interval control, and asynchronous processing remain necessary.

\subsection{Limitations}

Several limitations should be considered when interpreting the results. First, the benchmarking was carried out on a single target platform, the Jetson Orin Nano Dev Kit, so the absolute latency and hardware measurements should not be assumed to transfer directly to other embedded devices. Second, the model training and evaluation were based on a selected set of football broadcast datasets, which means that generalisation to lower-quality footage, unusual camera angles, severe occlusion, or non-broadcast match settings remains uncertain. Third, the full-stack analysis focused on a constrained set of match statistics and replay conditions rather than a fully production-ready analytics system. Finally, while a broad set of model artefacts was evaluated, the study necessarily sampled only part of the wider compression design space, especially for combined techniques and longer-duration deployment behaviour.

Despite these limitations, the thesis makes a useful contribution by linking model compression analysis to an operational soccer analytics pipeline rather than treating benchmark performance in isolation. The results show not only which techniques improve latency, but which ones remain meaningful once power, thermal behaviour, and end-to-end analytical usefulness are considered together.

\subsection{Conclusion}

This thesis investigated how computer vision models for soccer analysis can be trained, compressed, and deployed efficiently on embedded hardware. Across 488 trained and benchmarked model artefacts, the results showed that deployment-focused optimisation can substantially improve the practicality of edge-based sports analytics while preserving useful analytical performance. Hardware acceleration provided the safest and most broadly effective first step, quantisation delivered the strongest reductions in latency and operating cost, pruning proved more suitable for object detection than for pitch keypoint localisation, and knowledge distillation was most valuable as an accuracy-preserving support technique rather than as a direct speed-up method.

The work also demonstrated that model-level improvements can be carried through to a full-stack soccer analysis system for detection, localisation, possession estimation, and pass detection on the Jetson Orin Nano. Although the final system did not fully meet the most ambitious throughput target, it achieved a practically useful operating point at 13.02 FPS with moderate GPU utilisation and meaningful analytical output. Overall, the thesis shows that real-time or near-real-time soccer analytics on constrained edge hardware is achievable, but that the best deployment candidates are balanced compressed models rather than the most aggressively optimised ones. This provides a clear foundation for future work on system-level optimisation and broader match understanding, and supports the wider goal of low-cost, on-site analytics for teams, broadcasters, and sports science applications.
